\documentclass[a4paper,dvipdfmx]{jsarticle}
\usepackage[width=18cm,height=26cm]{geometry}
\usepackage{graphicx}
\usepackage{url}
\usepackage{listings,jlisting}
\lstset{
  basicstyle={\ttfamily},
  identifierstyle={\small},
  commentstyle={\smallitshape},
  keywordstyle={\small\bfseries},
  ndkeywordstyle={\small},
  stringstyle={\small\ttfamily},
  frame={tb},
  breaklines=true,
  columns=[l]{fullflexible},
  numbers=left,
  xrightmargin=0zw,
  xleftmargin=3zw,
  numberstyle={\scriptsize},
  stepnumber=1,
  numbersep=1zw,
  lineskip=-0.5ex
}
\begin{document}
\author{2345678901 名無しの権兵衛}
\title{情報科学概論第2回課題}
\maketitle
\section*{課題1}

\section*{課題2}

\section*{課題3}
\begin{center}
\IfFileExists{sha256.png}
{\includegraphics[width=0.7\linewidth]{sha256.png}}
{\framebox[0.7\linewidth]{\raisebox{-2.5cm}{\rule{0pt}{5cm}}スクリーンショットsha256.png}}
\end{center}

\section*{課題4}
\begin{enumerate}
\item 
\begin{center}
表1 2つの素数の積の素因数分解にかかる時間\\
\begin{tabular}{r|r}
\hline
\hline
\multicolumn{1}{c|}{各素数のビット数} &
\multicolumn{1}{c}{計算時間 [ms]} \\
\hline
% 以下の & と \\の間に結果を書き込むこと
50 &       \\
60 &       \\
70 &       \\
80 &       \\
90 &       \\
\hline
\end{tabular}
\end{center}
\item
\begin{description}
\item[【パスワード】]

\item[【平文】]

\item[【実行に要した時間】]

\end{description}
\end{enumerate}

\section*{ミニッツペーパー}
今回の授業の感想、疑問点などを書いてください。

\section*{謝辞}

\section*{参考文献}
% URLを記すときには、\url{https://www.ist.aichi-pu.ac.jp/}のように\urlコマンドを使うこと。

\end{document}
